\documentclass[11pt,a4paper]{article}

\usepackage[a4paper,top=18mm,bottom=20mm,left=17mm,right=17mm,headheight=15pt]{geometry}
\usepackage{fontspec}
\usepackage{xcolor}
\usepackage{array}
\usepackage{tabularx}
\usepackage{enumitem}
\usepackage{fancyhdr}
\usepackage{hyperref}
\usepackage{needspace}

\setmainfont{Arial}
\setsansfont{Arial}

\definecolor{Primary}{HTML}{235347}
\definecolor{Accent}{HTML}{E98B2A}
\definecolor{Soft}{HTML}{EEF4F1}
\definecolor{LineGray}{HTML}{B7C3BE}
\definecolor{TextGray}{HTML}{4E5A56}

\newcommand{\documentlabel}{Angket Terbuka Pelanggan GrabFood Nasi Gerilya}

\hypersetup{
  colorlinks=true,
  linkcolor=Primary,
  urlcolor=Primary,
  pdftitle={Angket Terbuka Pelanggan GrabFood Nasi Gerilya},
  pdfauthor={Instrumen Penelitian Skripsi}
}

\pagestyle{fancy}
\fancyhf{}
\fancyhead[L]{\small\textcolor{Primary}{\documentlabel}}
\fancyhead[R]{\small\textcolor{TextGray}{Instrumen Penelitian}}
\fancyfoot[C]{\small\textcolor{TextGray}{Halaman \thepage}}
\renewcommand{\headrulewidth}{0.4pt}
\renewcommand{\footrulewidth}{0pt}

\setlength{\parindent}{0pt}
\setlength{\parskip}{5pt}
\renewcommand{\arraystretch}{1.25}
\hyphenpenalty=10000
\exhyphenpenalty=10000
\emergencystretch=2em
\sloppy
\newcolumntype{Y}{>{\raggedright\arraybackslash}X}

\newcommand{\sectionbar}[1]{%
  \Needspace{4\baselineskip}
  \vspace{5mm}
  \noindent\colorbox{Soft}{%
    \parbox{\dimexpr\linewidth-2\fboxsep\relax}{%
      \large\bfseries\textcolor{Primary}{#1}
    }%
  }\par\vspace{2mm}
}

\newcommand{\answerlines}{%
  \par\vspace{1.2mm}
  \noindent\textcolor{LineGray}{\rule{\linewidth}{0.35pt}}\par\vspace{4.8mm}
  \noindent\textcolor{LineGray}{\rule{\linewidth}{0.35pt}}\par\vspace{3.2mm}
}

\newcommand{\question}[1]{%
  \Needspace{6\baselineskip}
  \item #1
  \answerlines
}

\newcommand{\widefield}[1]{%
  \Needspace{4\baselineskip}
  \textbf{#1}\par
  \noindent\textcolor{LineGray}{\rule{\linewidth}{0.35pt}}\par\vspace{4mm}
}

\begin{document}

\begin{titlepage}
\thispagestyle{empty}
\vspace*{10mm}
{\Huge\bfseries\textcolor{Primary}{Angket Terbuka}}\par
{\Huge\bfseries\textcolor{Primary}{Pelanggan GrabFood}}\par
{\Huge\bfseries\textcolor{Primary}{Nasi Gerilya}}\par
\vspace{3mm}
{\Large\textcolor{Accent}{Instrumen Pengumpulan Data Wawancara Pelanggan}}\par
\vspace{7mm}
\textcolor{Accent}{\rule{0.92\linewidth}{1.4pt}}\par
\vspace{10mm}

\begin{tabularx}{\linewidth}{>{\bfseries\color{Primary}}p{37mm} Y}
Topik penelitian & Pengalaman, persepsi, kepuasan, dan loyalitas pelanggan GrabFood terhadap Nasi Gerilya.\\
Sasaran informan & Pelanggan yang pernah membeli Nasi Gerilya melalui aplikasi GrabFood.\\
Metode pengisian & Pertanyaan terbuka. Jawaban dapat ditulis langsung atau dicatat oleh pewawancara saat interview.\\
Tujuan & Menggali alasan pembelian, penilaian terhadap produk, harga, pelayanan, promosi, kepuasan, pembelian ulang, dan saran perbaikan.\\
Kerahasiaan & Identitas informan dijaga kerahasiaannya dan data digunakan hanya untuk keperluan penelitian.\\
Perkiraan durasi & 45 sampai 60 menit, menyesuaikan kedalaman jawaban informan.\\
\end{tabularx}

\vfill
\colorbox{Soft}{%
  \parbox{\dimexpr\linewidth-2\fboxsep\relax}{%
    \raggedright
    \textbf{\textcolor{Primary}{Petunjuk singkat:}}
    Jawablah setiap pertanyaan berdasarkan pengalaman pribadi saat membeli Nasi Gerilya melalui GrabFood. Tidak ada jawaban benar atau salah. Jelaskan alasan, contoh, atau pengalaman nyata jika memungkinkan.
  }%
}
\vspace*{8mm}
\end{titlepage}

\sectionbar{A. Data Pelaksanaan Wawancara}
\widefield{Kode informan}
\widefield{Tanggal wawancara}
\widefield{Tempat atau media wawancara}
\widefield{Nama pewawancara}
\widefield{Kesediaan informan mengikuti wawancara}
\widefield{Catatan awal pewawancara}

\sectionbar{B. Identitas Informan}
\begin{enumerate}[label=\arabic*.,leftmargin=*,itemsep=2mm]
\question{Boleh ceritakan nama atau inisial Anda dan usia Anda?}
\question{Apa pekerjaan atau aktivitas utama Anda saat ini?}
\question{Di daerah mana Anda tinggal atau biasa beraktivitas?}
\question{Bagaimana kebiasaan Anda dalam membeli makanan secara online?}
\end{enumerate}

\sectionbar{C. Pengalaman Menggunakan GrabFood}
\begin{enumerate}[label=\arabic*.,leftmargin=*,itemsep=2mm,resume]
\question{Sejak kapan Anda mulai menggunakan GrabFood?}
\question{Seberapa sering Anda biasanya memesan makanan melalui GrabFood?}
\question{Dalam situasi seperti apa Anda biasanya menggunakan GrabFood untuk membeli makanan?}
\question{Jenis makanan seperti apa yang paling sering Anda pesan melalui GrabFood?}
\question{Apa saja hal yang biasanya Anda pertimbangkan sebelum memilih makanan di GrabFood?}
\question{Bagaimana pengaruh harga, promo, rating, foto, dan jarak restoran terhadap keputusan Anda saat memesan makanan?}
\end{enumerate}

\sectionbar{D. Pengalaman Mengenal Nasi Gerilya}
\begin{enumerate}[label=\arabic*.,leftmargin=*,itemsep=2mm,resume]
\question{Ceritakan bagaimana awalnya Anda mengetahui Nasi Gerilya.}
\question{Apa yang membuat Anda tertarik untuk mencoba Nasi Gerilya melalui GrabFood?}
\question{Bagaimana kesan pertama Anda saat melihat Nasi Gerilya di aplikasi GrabFood?}
\question{Apa menu Nasi Gerilya yang pernah atau paling sering Anda pesan?}
\question{Mengapa Anda memilih menu tersebut?}
\question{Bagaimana pengalaman pertama Anda saat membeli Nasi Gerilya melalui GrabFood?}
\end{enumerate}

\sectionbar{E. Proses Keputusan Pembelian}
\begin{enumerate}[label=\arabic*.,leftmargin=*,itemsep=2mm,resume]
\question{Ceritakan proses Anda saat memutuskan membeli Nasi Gerilya di GrabFood.}
\question{Apa faktor utama yang membuat Anda akhirnya memilih Nasi Gerilya dibanding pilihan makanan lain?}
\question{Bagaimana peran foto menu dalam memengaruhi keputusan Anda?}
\question{Bagaimana peran nama menu dan deskripsi menu dalam memengaruhi keputusan Anda?}
\question{Bagaimana peran rating atau ulasan pelanggan lain dalam keputusan Anda membeli?}
\question{Bagaimana pengaruh promo atau diskon terhadap minat Anda membeli Nasi Gerilya?}
\question{Apa yang biasanya membuat Anda ragu sebelum membeli makanan melalui GrabFood?}
\question{Hal apa yang membuat Anda yakin untuk membeli Nasi Gerilya?}
\end{enumerate}

\sectionbar{F. Penilaian terhadap Produk}
\begin{enumerate}[label=\arabic*.,leftmargin=*,itemsep=2mm,resume]
\question{Bagaimana pendapat Anda tentang rasa makanan Nasi Gerilya?}
\question{Bagian rasa apa yang paling Anda sukai dari Nasi Gerilya?}
\question{Bagian rasa apa yang menurut Anda masih perlu diperbaiki?}
\question{Bagaimana pendapat Anda tentang porsi makanan Nasi Gerilya?}
\question{Bagaimana kesesuaian porsi dengan kebutuhan atau harapan Anda?}
\question{Bagaimana pendapat Anda tentang tampilan makanan saat diterima?}
\question{Bagaimana kondisi makanan saat sampai di tangan Anda?}
\question{Bagaimana pendapat Anda tentang suhu makanan saat diterima?}
\question{Bagaimana pendapat Anda tentang kualitas bahan makanan yang digunakan?}
\question{Bagaimana pendapat Anda tentang variasi menu Nasi Gerilya di GrabFood?}
\question{Menu seperti apa yang menurut Anda perlu ditambahkan?}
\question{Bagaimana kesesuaian antara foto menu di GrabFood dengan makanan yang Anda terima?}
\question{Bagaimana kesesuaian antara deskripsi menu dengan makanan yang Anda terima?}
\question{Bagaimana pendapat Anda tentang tingkat kepedasan, bumbu, atau cita rasa makanan?}
\question{Apa yang membedakan rasa Nasi Gerilya dari makanan sejenis lainnya?}
\end{enumerate}

\sectionbar{G. Penilaian terhadap Kemasan}
\begin{enumerate}[label=\arabic*.,leftmargin=*,itemsep=2mm,resume]
\question{Bagaimana pendapat Anda tentang kemasan makanan Nasi Gerilya?}
\question{Bagaimana kondisi kemasan saat makanan sampai?}
\question{Bagaimana kemampuan kemasan dalam menjaga makanan agar tidak tumpah atau rusak?}
\question{Bagaimana kesan Anda terhadap kebersihan makanan dan kemasan?}
\question{Apa yang menurut Anda perlu diperbaiki dari sisi kemasan?}
\question{Kemasan seperti apa yang menurut Anda lebih menarik atau lebih nyaman digunakan?}
\end{enumerate}

\sectionbar{H. Penilaian terhadap Harga}
\begin{enumerate}[label=\arabic*.,leftmargin=*,itemsep=2mm,resume]
\question{Bagaimana pendapat Anda tentang harga Nasi Gerilya di GrabFood?}
\question{Bagaimana kesesuaian harga dengan rasa makanan yang Anda dapatkan?}
\question{Bagaimana kesesuaian harga dengan porsi makanan?}
\question{Bagaimana kesesuaian harga dengan kualitas makanan secara keseluruhan?}
\question{Bagaimana pendapat Anda tentang biaya tambahan seperti ongkir dan biaya layanan saat membeli Nasi Gerilya?}
\question{Dalam kondisi tanpa promo, bagaimana minat Anda untuk tetap membeli Nasi Gerilya?}
\question{Bagaimana perbandingan harga Nasi Gerilya dengan makanan sejenis di GrabFood?}
\question{Menurut Anda, harga seperti apa yang paling sesuai untuk produk Nasi Gerilya?}
\end{enumerate}

\sectionbar{I. Penilaian terhadap Pelayanan melalui GrabFood}
\begin{enumerate}[label=\arabic*.,leftmargin=*,itemsep=2mm,resume]
\question{Bagaimana pengalaman Anda saat memesan Nasi Gerilya melalui GrabFood?}
\question{Bagaimana pendapat Anda tentang kecepatan restoran dalam memproses pesanan?}
\question{Bagaimana pendapat Anda tentang waktu tunggu dari pemesanan sampai makanan diterima?}
\question{Ceritakan pengalaman Anda terkait kesesuaian pesanan yang diterima.}
\question{Bagaimana pengalaman Anda jika memberikan catatan khusus pada pesanan?}
\question{Bagaimana restoran memenuhi permintaan atau catatan khusus dari pelanggan?}
\question{Kendala apa yang pernah Anda alami saat memesan Nasi Gerilya melalui GrabFood?}
\question{Bagaimana penyelesaian kendala tersebut menurut pengalaman Anda?}
\question{Bagaimana peran driver GrabFood dalam pengalaman Anda menerima pesanan Nasi Gerilya?}
\end{enumerate}

\sectionbar{J. Tampilan dan Informasi Nasi Gerilya di GrabFood}
\begin{enumerate}[label=\arabic*.,leftmargin=*,itemsep=2mm,resume]
\question{Bagaimana pendapat Anda tentang tampilan Nasi Gerilya di aplikasi GrabFood?}
\question{Bagaimana kemudahan Anda dalam menemukan Nasi Gerilya di GrabFood?}
\question{Bagaimana kejelasan informasi menu yang tersedia di aplikasi?}
\question{Bagaimana kejelasan informasi harga di aplikasi?}
\question{Bagaimana pendapat Anda tentang foto produk yang ditampilkan?}
\question{Bagaimana pendapat Anda tentang nama-nama menu yang digunakan?}
\question{Informasi apa yang menurut Anda masih kurang di halaman GrabFood Nasi Gerilya?}
\question{Tampilan seperti apa yang menurut Anda bisa membuat Nasi Gerilya lebih menarik di GrabFood?}
\end{enumerate}

\sectionbar{K. Promosi}
\begin{enumerate}[label=\arabic*.,leftmargin=*,itemsep=2mm,resume]
\question{Bagaimana Anda mengetahui promo Nasi Gerilya di GrabFood?}
\question{Bagaimana pengaruh promo terhadap keputusan Anda membeli Nasi Gerilya?}
\question{Promo seperti apa yang paling menarik bagi Anda?}
\question{Bagaimana pendapat Anda tentang promo yang pernah ditawarkan Nasi Gerilya?}
\question{Selain diskon harga, bentuk promosi apa yang menurut Anda bisa menarik pelanggan?}
\question{Bagaimana media sosial, rekomendasi teman, atau ulasan pelanggan memengaruhi minat Anda terhadap Nasi Gerilya?}
\end{enumerate}

\sectionbar{L. Kepuasan Pelanggan}
\begin{enumerate}[label=\arabic*.,leftmargin=*,itemsep=2mm,resume]
\question{Bagaimana tingkat kepuasan Anda setelah membeli Nasi Gerilya melalui GrabFood?}
\question{Hal apa yang membuat Anda puas saat membeli Nasi Gerilya?}
\question{Hal apa yang membuat Anda kurang puas atau kecewa?}
\question{Bagaimana kesesuaian pengalaman membeli Nasi Gerilya dengan harapan Anda sebelumnya?}
\question{Bagaimana nilai yang Anda rasakan antara uang yang dikeluarkan dan makanan yang diterima?}
\question{Ceritakan pengalaman paling berkesan saat membeli Nasi Gerilya.}
\question{Ceritakan pengalaman yang kurang menyenangkan, jika ada, saat membeli Nasi Gerilya.}
\end{enumerate}

\sectionbar{M. Loyalitas dan Pembelian Ulang}
\begin{enumerate}[label=\arabic*.,leftmargin=*,itemsep=2mm,resume]
\question{Apa yang membuat Anda ingin membeli kembali Nasi Gerilya?}
\question{Apa yang bisa membuat Anda tidak ingin membeli kembali Nasi Gerilya?}
\question{Dalam kondisi seperti apa Anda biasanya terpikir untuk membeli Nasi Gerilya lagi?}
\question{Bagaimana kemungkinan Anda merekomendasikan Nasi Gerilya kepada teman atau keluarga?}
\question{Apa alasan Anda merekomendasikan atau tidak merekomendasikan Nasi Gerilya?}
\question{Hal apa yang perlu dipertahankan agar Anda tetap menjadi pelanggan Nasi Gerilya?}
\question{Hal apa yang perlu ditingkatkan agar Anda lebih sering membeli Nasi Gerilya?}
\end{enumerate}

\sectionbar{N. Perbandingan dengan Kompetitor}
\begin{enumerate}[label=\arabic*.,leftmargin=*,itemsep=2mm,resume]
\question{Makanan atau restoran apa yang biasanya Anda bandingkan dengan Nasi Gerilya?}
\question{Bagaimana perbandingan Nasi Gerilya dengan makanan sejenis dari segi rasa?}
\question{Bagaimana perbandingan Nasi Gerilya dengan makanan sejenis dari segi harga?}
\question{Bagaimana perbandingan Nasi Gerilya dengan makanan sejenis dari segi porsi?}
\question{Bagaimana perbandingan Nasi Gerilya dengan makanan sejenis dari segi kemasan?}
\question{Bagaimana perbandingan Nasi Gerilya dengan makanan sejenis dari segi pelayanan?}
\question{Apa keunggulan utama Nasi Gerilya dibandingkan pesaingnya?}
\question{Apa kelemahan utama Nasi Gerilya dibandingkan pesaingnya?}
\end{enumerate}

\sectionbar{O. Masukan dan Harapan}
\begin{enumerate}[label=\arabic*.,leftmargin=*,itemsep=2mm,resume]
\question{Apa hal paling Anda sukai dari Nasi Gerilya?}
\question{Apa hal yang paling perlu diperbaiki oleh Nasi Gerilya?}
\question{Menu baru seperti apa yang menurut Anda cocok untuk Nasi Gerilya?}
\question{Bagaimana saran Anda untuk meningkatkan kualitas rasa Nasi Gerilya?}
\question{Bagaimana saran Anda untuk meningkatkan tampilan dan kemasan produk?}
\question{Bagaimana saran Anda untuk meningkatkan pelayanan Nasi Gerilya di GrabFood?}
\question{Bagaimana saran Anda agar Nasi Gerilya lebih menarik bagi pelanggan baru?}
\question{Apa harapan Anda terhadap Nasi Gerilya ke depannya?}
\question{Apakah ada hal lain yang ingin Anda sampaikan tentang pengalaman membeli Nasi Gerilya melalui GrabFood?}
\end{enumerate}

\sectionbar{P. Catatan Tambahan Pewawancara}
\widefield{Ringkasan jawaban penting}
\widefield{Kutipan langsung yang relevan}
\widefield{Kesan nonverbal atau situasi wawancara}
\widefield{Kode tema sementara}

\end{document}
